% YAAC Another Awesome CV LaTeX Template
%
% This template has been downloaded from:
% https://github.com/darwiin/yaac-another-awesome-cv
%
% Author:
% Christophe Roger
%
% Template license:
% CC BY-SA 4.0 (https://creativecommons.org/licenses/by-sa/4.0/)
%Section: Work Experience at the top
\sectionTitle{Experiência}{\faSuitcase}
%\renewcommand{\labelitemi}{$\bullet$}
\begin{experiences}
    \experience
        {Setembro 2026} {H\&Fti}{Arquiteto de Software}{}
        {Presente} { Condução e desenvolvimento de projetos de software sob medida, com foco em arquiteturas distribuídas, entrega rápida de MVPs e engenharia assistida por inteligência artificial.\\
        Principais frentes e entregas:
            \begin{itemize}
               \item Desenvolvimento de Aplicações \& MVPs: Arquitetura e construção de soluções completas de ponta a ponta (stacks modernas JAPS).
               \item Estratégia de Dados \& Nuvem: Modelagem e escolha estratégica (Relacional vs. NoSQL) para escalabilidade.
               \item DevOps \& Infraestrutura: Implementação de pipelines automatizadas (CI/CD), orquestração de containers (Docker/Rancher) e integração com edge/CDN.
               \item Aceleração com IA: Integração de fluxos de desenvolvimento com IA (vibecoding contextualizado, análise preditiva de bugs e documentação automática).
              \end{itemize}
            }
            {Arquitetura Distribuída, IA, CI/CD, Docker, Rancher, Cloudflare, Relacional/NoSQL, Vibecoding}
    \emptySeparator
    \experience
        {Agosto 2025} {nstech}{Líder de equipe sênior BPORT}{}
        {Setembro 2026} { Atuação na liderança técnica e de equipe do projeto BPORT, com foco em maturidade de engenharia, eficiência de processos e pioneirismo na adoção de Inteligência Artificial aplicada ao ciclo de desenvolvimento.
            \begin{itemize}
               \item Inovação com IA na Engenharia: Idealização e condução de soluções automatizadas para documentação e mapeamento preditivo.
               \item Engenharia de Contexto \& Metodologia: Estruturação de regras e contextos avançados para desenvolvimento assistido por IA (vibecoding).
               \item Transformação Cultural: Difusão e capacitação do uso de ferramentas de IA entre os times.
               \item Evolução do SDLC: Refinamento contínuo do Git Flow e otimização de todo o fluxo de desenvolvimento.
              \end{itemize}
            }
            {Liderança Técnica, Inteligência Artificial, SDLC, Git Flow, Vibecoding}
    \emptySeparator
    \experience
        {Janeiro 2025} {PETROBRAS}{Líder de equipe técnica}{}
        {Julho 2025} { Liderança técnica da equipe responsável por soluções de alta criticidade e escala para o ecossistema Petrobras no setor logístico/portuário.
            \begin{itemize}
               \item Liderança Técnica \& Ritos: Condução do time com metodologias ágeis (Scrum/Kanban) e cultura de code review.
               \item Padronização de Versionamento: Reestruturação e refinamento do Git Flow da equipe.
               \item Arquitetura e Sustentação: Acompanhamento de entregas críticas com foco em resiliência e escalabilidade nas integrações.
              \end{itemize}
            }
            {Scrum, Kanban, Git Flow, Arquitetura, Liderança Técnica}
    \emptySeparator
    \experience
        {Outubro 2024} {nstech}{Desenvolvedor Full Stack | GRAILS}{}
        {Dezembro 2024} { Atuação direta no desenvolvimento e evolução de aplicações robustas e escaláveis voltadas ao setor portuário e logístico.
            \begin{itemize}
               \item Desenvolvimento Backend \& Fullstack: Funcionalidades com foco em performance e manutenibilidade.
               \item DevOps \& Integração Contínua: Gerenciamento de código com GitLab e automação de deploys com Jenkins.
               \item Conteinerização: Hospedagem e orquestração de aplicações via Rancher e Docker.
               \item Participação ativa em todas as fases do ciclo de desenvolvimento, da análise à produção.
              \end{itemize}
            }
            {Grails, Spring Boot, Spring Data JPA, GitLab, Jenkins, Docker, Rancher}
    \emptySeparator
    \experience
        {Fevereiro 2024} {H\&Fti}{Desenvolvedor FullStack}{}
        {Setembro 2024} { Analisar fluxos de processos, criar projetos de desenvolvimento e desenvolver soluções compatíveis com as necessidades setoriais. Consultoria de tecnologia e análise de sistemas.
            \begin{itemize}
               \item Desenvolvimento utilizando Java com Angular.
               \item Modelagem e implementação em bancos de dados SQL (PostgreSQL) e NoSQL (MongoDB).
               \item Prestação de suportes a bugs em sistemas legados em RubyOnRails.
              \end{itemize}
            }
            {Java, Angular, PostgreSQL, MongoDB, RubyOnRails}
    \emptySeparator
    \experience
        {Outubro 2022} {Capgemini}{Team Lead}{}
        {Fevereiro 2024} { Liderar o desenvolvimento da equipe TISS, desenvolver estratégias de desenvolvimento e multiplicar conhecimento do projeto. Alinhar fluxo de desenvolvimento com Scrum Master e PO. Manter comunicação com equipes externa dentro das necessidades do produto TISS.
            \begin{itemize}
               \item Gestão técnica e liderança de equipe.
               \item Comunicação constante entre times e alinhamento ágil.
               \item Desenvolvimento de Sistemas com JAVA Web, Angular, Plataformas integradas com JSON e XML.
              \end{itemize}
            }
            {Liderança, JAVA, WebSphere, Angular, XML, JSON, CI/CD}
    \emptySeparator
    \experience
        {Abril 2022}   {Capgemini}{Desenvolvedor | Back-end | Java}{}
        {Setembro 2022} { Designado para desenvolver em sistema Java que recebia uma BOOK Cobol tratava e retornava enviando um XML para um servidor em alta plataforma.
            \begin{itemize}
                \item Montagem de arquivos Java XML dentro dos padrões da ANVISA.
                \item Proposição e desenvolvimento de nova solução JSON para ganho de performance, evitando estouro de memória.
                \item Utilização do ambiente WebSphere para gerenciar libs, serviços de banco e web services.
            \end{itemize}
            }
            {JAVA, WebSphere, Angular, XML, Jenkins, JSON, CI/CD, Struts, JSF, SOAP}
    \emptySeparator
    \experience
        {Fevereiro 2022}   {Blue Technology}{Desenvolvedor | Back-end | Java}{}
        {Agosto 2022} { Atuando como desenvolvedor de base Java desenvolvi em soluções com Kotlin, ReactJS, Angular, AngularJS, Groovy e Java. Em microserviços utilizando CI/CD dentro da ferramenta GitLab.
                    \begin{itemize}
                        \item Ambiente de Micro Serviços.
                        \item Marcações para Observabilidade e Métricas (Actuator, Grafana, Prometheus).
                        \item Desenvolvimento em TDD e mensageria com Kafka.
                        \item Integração e Deploy Contínuos; testes automatizados (Cypress).
                    \end{itemize}
                    }
                    {JAVA, Kotlin, Angular, Groovy, TDD, Cypress, Gitlab, Jenkins, Kafka, Prometheus}
    \emptySeparator
    \experience
            {Setembro 2021}   {Pitang}{Desenvolvedor | Back-end | Java}{}
            {Dezembro 2021} { Contratado para atuar desenvolvendo um sistema de Gestão e Avaliação de Projetos Acadêmicos para empresas de Petróleo.
                    \begin{itemize}
                        \item JAVA com Spring boot, data e batch;
                        \item Banco de dados MongoDB;
                        \item Utilização de Kanban e Scrum;
                        \item Rest API e mapeamento de endpoints.
                    \end{itemize}
                    }
                    {JAVA, Spring Boot, Spring Data, Spring Batch, RESTful, MongoDB, GitLab}
    \emptySeparator
    \experience
            {Julho 2017} {Corpo de Bombeiros Militar de Pernambuco}{Engenheiro de Software | Full Stack}{}
            {Dezembro 2019} { CTIC – Centro de Tecnologia da Informação e Comunicação\\
            Desafio de criar soluções para setores específicos demandados dentro da Corporação em RubyOnRails e Java.
                    \begin{itemize}
                       \item Projeto entregue em 18 meses com 2 produtos diferentes para confecção de carteiras de identidades (biblioteca PRAWN).
                       \item Desenvolvimento com Spring Boot, MongoDB e ReactJS para mineração e controle num Dashboard em tempo real.
                       \item Ministrava aulas sobre tecnologias da Corporação, LGPD e Ética Profissional.
                      \end{itemize}
                    }
                    {Java, Spring boot, Angular, RubyOnRails, ReactJS, PostgreSQL, MongoDB, Docker, Prawn}
  \emptySeparator
  \experience
        {Outubro 2014} {MV S/A}{Engenheiro de Software | Back-end | I18N}{}
        {Fevereiro 2015} { Empresa Multinacional baseada em framework próprio. Captura de problemas no suporte ao I18N, tornando-me responsável do suporte total à I18N em todos os setores e partes da fábrica.
                    \begin{itemize}
                        \item Soluções de internacionalização com i18n.
                        \item Integração entre equipes, desonerando o tempo de outras áreas e aumentando reuso.
                    \end{itemize}
                    }
                    {JAVA, i18n}
  \emptySeparator
 \experience
        {Maio 2013}   {Procenge }{Desenvolvedor | Back-end | Java}{}
        {Julho 2014} { Atuando no desenvolvimento de uma plataforma de gerência de recursos naturais (gás).
                \begin{itemize}
                    \item Desenvolvimento com Java Spring MVC com JSF/JSP, Struts.
                    \item Comunicação com banco através de Hibernate e Subversion.
                    \item Utilizando ferramentas: IReport, NetBeans, MySQL, Eclipse, Oracle, Toad, StarUML, HTML, CSS, JS, JQuery, Ajax.
                    \end{itemize}
                }
                {Eclipse, JAVA, Spring MVC, hibernate, MySQL, Oracle, SubVersion }
  \emptySeparator
  \experience
        {Novembro 2012} {SoftexRecife}{Engenheiro de Software | Desenvolvedor}{}
        {Abril 2013}{ Programa de Residência em Software da SOFTEX Recife, ministrado pela PROCENGE.
                        \begin{itemize}
                            \item Módulos: SGP, C\# Básico, Lógica de Programação, SQL ANSI, Crystal Report, Dev Web com .Net, Centura.
                            \item Desenvolvimento com Java e .NET, MySQL e Treinamento na IDE Eclipse.
                        \end{itemize}
                    }
                    {Eclipse, JAVA, .NET, MySQL}
    \emptySeparator
\end{experiences}
